%% 第三章

\chapter{UAV-DualCog双认知测试基准设计}
\chaptermark{UAV-DualCog双认知测试基准设计}

\section{总体设计思想}

UAV-DualCog 的核心设计思想是从无人机具身推理需求出发，将多模态大模型能力评测从单一语义问答扩展为“认知维度、观测媒介、证据形式”相结合的系统评测。无人机的观察不是来自被动摄像机，而是由主动运动、飞行高度、航向变化和相机视角共同决定。因此，模型必须同时回答两个问题：一是无人机是否理解自身状态，二是无人机是否理解外部环境。

\begin{figure}[htbp]
    \centering
    \includegraphics[width=\textwidth]{fig_benchmark_overview.pdf}
    \caption{UAV-DualCog测试基准总体结构}
    \label{fig:benchmark-overview}
\end{figure}

如\autoref{fig:benchmark-overview}所示，UAV-DualCog 将自我状态认知和环境情况认知作为两条顶层能力主线，并通过图像任务和视频任务共同测量。图像任务强调离散观察下的空间推理和目标定位，视频任务强调连续观察下的行为理解、可见性判断和时间定位。两类媒介共同服务于同一个问题：模型是否能够在无人机主动观测场景中同时理解“自己”和“外界”，并用可检验的空间或时间证据支撑判断。

\section{双认知能力定义}

本文将无人机双认知能力定义为自我状态认知和环境情况认知两部分。对于无人机中心观测 \(o\)，其认知目标可以表示为
\begin{equation}
    \mathcal{C}(o)=\{c_s,c_e\},
    \label{eq:dual-cognition}
\end{equation}
其中，\(c_s\) 表示自我状态认知，\(c_e\) 表示环境情况认知。

自我状态认知关注无人机对自身状态的理解，包括无人机相对于地标的位置、运动导致的未来观察变化、当前执行的飞行行为以及行为发生的时间区间。环境情况认知关注无人机对外部世界的理解，包括地标相对于无人机前向的位置、根据地标位置应采取的动作方向、目标地标在视频中的可见次数和可见时间区间。

这一区分具有实际意义。对于普通视觉问答任务，模型可以只描述图像中有什么目标或目标之间存在什么关系；对于无人机而言，目标外观和空间关系都受到无人机自身位置、朝向和运动的影响。如果模型不能理解自身状态，那么它对外部环境的判断也可能缺乏稳定参照系。因此，双认知能力不是两个互不相关的能力集合，而是无人机具身推理中相互条件化、相互约束的两个方面。

\section{六类任务体系}

在双认知框架基础上，本文设计了四类图像任务和两类视频任务。六类任务的具体设计如\autoref{tab:task-system}所示。

\begin{table}[htbp]
    \centering
    \caption{UAV-DualCog六类任务体系}
    \label{tab:task-system}
    \begin{tabularx}{\textwidth}{p{3.0cm}p{1.6cm}p{2.0cm}X}
        \toprule
        任务名称 & 模态 & 认知维度 & 任务定义 \\
        \midrule
        地标相对位置推理 & 图像 & 自我状态认知 & 给定地标中心参考视图和当前第一视角观测，判断无人机相对于目标地标的位置，并输出当前图像中的目标框。 \\
        未来观察预测 & 图像 & 自我状态认知 & 给定参考视图和候选未来观察图像，根据动作变化描述选择正确未来观察，并输出对应目标框。 \\
        自我参照位置推理 & 图像 & 环境情况认知 & 在当前第一视角观测中判断目标地标相对于无人机前向的方位，并输出目标框。 \\
        地标驱动动作决策 & 图像 & 环境情况认知 & 根据当前观测中的目标地标位置判断无人机应朝哪个方向移动，并输出目标框。 \\
        飞行行为识别与时间定位 & 视频 & 自我状态认知 & 从第一视角飞行视频中识别无人机正在执行的飞行行为，并给出对应时间区间。 \\
        地标可见次数与区间推理 & 视频 & 环境情况认知 & 判断目标地标在飞行视频中的出现次数，并给出全部可见时间区间。 \\
        \bottomrule
    \end{tabularx}
\end{table}

六类任务在结构上保持统一。图像任务要求输出离散选项和标准化二维目标框；视频任务要求输出行为或可见次数，并给出秒级时间区间。这样设计的原因在于，UAV-DualCog 不希望只评估模型是否能够给出看似正确的自然语言解释，而是要求模型在严格结构化输出下同时完成语义判断和证据定位。

\section{输出契约与评测指标}

UAV-DualCog 为图像和视频任务分别设计了严格的结构化输出契约。图像任务统一使用 \texttt{answer\_option\_id} 和 \texttt{bbox\_xyxy\_norm} 字段，视频任务则使用行为答案列表、时间区间或可见次数相关字段。主要评测指标如\autoref{tab:evaluation-protocol}所示。

\begin{table}[htbp]
    \centering
    \caption{任务输出契约与主要评测指标}
    \label{tab:evaluation-protocol}
    \small
    \setlength{\tabcolsep}{4pt}
    \begin{tabularx}{\textwidth}{@{}p{3.0cm}p{2.2cm}p{3.5cm}X@{}}
        \toprule
        任务类型 & 认知维度 & 结构化输出 & 主要指标 \\
        \midrule
        图像任务 & \makecell[l]{自我状态认知\\或环境情况认知} & \makecell[l]{\texttt{answer\_option\_id}\\\texttt{bbox\_xyxy\_norm}} & \makecell[l]{选项准确率\\BBox Acc@50IoU\\BBox Mean IoU} \\
        复合行为识别 & 自我状态认知 & \makecell[l]{\texttt{answers[].option\_id}\\\texttt{answers[].intervals\_sec}} & \makecell[l]{Semantic Accuracy\\Temporal F1@0.5\\Mean tIoU} \\
        原子行为识别 & 自我状态认知 & \makecell[l]{\texttt{answers[].option\_id}\\\texttt{answers[].intervals\_sec}} & \makecell[l]{Semantic Accuracy\\Temporal F1@0.5\\Mean tIoU} \\
        地标可见性推理 & 环境情况认知 & \makecell[l]{\texttt{visible\_count}\\\texttt{visible\_intervals\_sec}} & \makecell[l]{Count Accuracy\\Segment F1@0.5\\Mean tIoU} \\
        \bottomrule
    \end{tabularx}
\end{table}

从评测协议可以看出，本文将目标框和时间区间作为正式指标，而不是附属可视化信息。这使得实验可以进一步分析模型是否存在“答对但证据错误”的现象。例如，图像任务中模型可能选对方向选项，却无法定位地标；视频任务中模型可能判断出目标出现次数，却无法给出完整可见区间。UAV-DualCog 通过将这些证据纳入指标，能够更真实地衡量无人机具身推理的可靠性。

\section{提示词协议设计}

为了保证不同模型在统一接口下可比较，本文对六类任务设计了固定提示词协议。每个任务由系统提示词、用户提示模板和 JSON 输出模式三部分组成。系统提示词规定无人机代理身份、输入类型、输出字段和禁止事项；用户提示词注入当前任务中的地标描述、参考图像、当前观测、候选选项或视频语境；JSON 输出模式则保证解析器可以稳定提取字段并计算指标。正文给出协议摘要，完整协议结构见附录\autoref{app:prompt-protocol}。

\begin{table}[htbp]
    \centering
    \caption{六类任务提示词设计摘要}
    \label{tab:prompt-summary}
    \begin{tabularx}{\textwidth}{p{3.1cm}X p{3.5cm}}
        \toprule
        任务名称 & 核心约束 & 输出契约 \\
        \midrule
        地标相对位置推理 & 指定无人机代理身份，输入参考视图和当前视图，必须同时给出选项与目标框。 & \texttt{\{answer, bbox\}} \\
        未来观察预测 & 从多个候选未来视图中选择正确观察，并给出选中图像中的目标框。 & \texttt{\{answer, bbox\}} \\
        自我参照位置推理 & 判断目标地标相对无人机前向方向，并在当前观测中定位目标。 & \texttt{\{answer, bbox\}} \\
        地标驱动动作决策 & 根据地标位置选择接近目标的飞行动作，并输出作为证据的目标框。 & \texttt{\{answer, bbox\}} \\
        飞行行为识别与时间定位 & 识别一个或多个飞行行为选项，并给出各自对应时间区间。 & \texttt{\{answers\}} \\
        地标可见次数与区间推理 & 输出目标地标可见次数和全部可见时间区间。 & \texttt{\{count, intervals\}} \\
        \bottomrule
    \end{tabularx}
\end{table}

这种协议层设计将答案字段和证据字段绑定在一起，避免模型通过宽泛自然语言描述绕开评测解析。对于后续能力优化而言，结构化输出也为监督微调和可验证奖励提供了基础。

\section{测试基准规模与统计}

UAV-DualCog 的完整资产池覆盖 18 个 AerialVLN 场景，总面积约为 5040232.00 m\(^2\)。经过自动构建和人工复核后，从 2626 个候选地标中保留 746 个有效地标，覆盖 8 个粗粒度类别和 166 个细粒度子类。当前发布的测试基准从中选择 12 个场景和 512 个有效地标，构建 4096 个图像任务样本和 2048 个视频任务样本。

图像子集包含四类任务，每类 1024 个样本，共使用 3826 张唯一源图像和 12288 次图像引用。视频子集包含两类任务，每类 1024 个样本，基于 1024 条飞行轨迹生成，总轨迹长度为 388661.79 m，总视频时长为 19633.30 s，约合 5 小时 27 分 13 秒。

\begin{figure}[htbp]
    \centering
    \includegraphics[width=\textwidth]{fig_benchmark_stat.pdf}
    \caption{UAV-DualCog测试基准统计信息}
    \label{fig:benchmark-stat}
\end{figure}

如\autoref{fig:benchmark-stat}所示，测试基准在任务数量、认知维度和媒介类型上保持均衡。每类任务均为 1024 个样本，可以避免某一任务因样本数量过大而主导总体评测结果。场景方面覆盖城市、乡镇、工业区、水滨区域、夜间场景等多种低空无人机常用环境，为模型能力分析提供了较丰富的场景基础。

\section{本章小结}

本章介绍了 UAV-DualCog 的总体设计、双认知定义、六类任务体系、输出契约、提示词协议和测试基准规模。该基准的关键特点在于将自我状态认知与环境情况认知同时纳入评测，并把空间目标框和时间区间作为正式证据指标。下一章将进一步介绍测试基准数据构建工具链的实现过程。
